\documentclass[10pt,a4paper]{amsart}
\usepackage{amsmath,amssymb,amsthm,graphicx}
\usepackage{verbatim}
\newenvironment{reference}{\comment}{\endcomment}
%\newenvironment{reference}{}{}
\newenvironment{slogan}{\comment}{\endcomment}
\newenvironment{history}{\comment}{\endcomment}

% For commutative diagrams we use Xy-pic
\usepackage[all]{xy}

% We use 2cell for 2-commutative diagrams.
\xyoption{2cell}
% Theorem environments.
%
\theoremstyle{plain}
\newtheorem{theorem}[subsection]{Theorem}
\newtheorem{proposition}[subsection]{Proposition}
\newtheorem{lemma}[subsection]{Lemma}

\theoremstyle{definition}
\newtheorem{definition}[subsection]{Definition}
\newtheorem{example}[subsection]{Example}
\newtheorem{exercise}[subsection]{Exercise}
\newtheorem{situation}[subsection]{Situation}

\theoremstyle{remark}
\newtheorem{remark}[subsection]{Remark}
\newtheorem{remarks}[subsection]{Remarks}

\numberwithin{equation}{subsection}

% Macros
%
\def\lim{\mathop{\mathrm{lim}}\nolimits}
\def\colim{\mathop{\mathrm{colim}}\nolimits}
\def\Spec{\mathop{\mathrm{Spec}}}
\def\Hom{\mathop{\mathrm{Hom}}\nolimits}
\def\Ext{\mathop{\mathrm{Ext}}\nolimits}
\def\SheafHom{\mathop{\mathcal{H}\!\mathit{om}}\nolimits}
\def\SheafExt{\mathop{\mathcal{E}\!\mathit{xt}}\nolimits}
\def\Sch{\mathit{Sch}}
\def\Mor{\mathop{\mathrm{Mor}}\nolimits}
\def\Ob{\mathop{\mathrm{Ob}}\nolimits}
\def\Sh{\mathop{\mathit{Sh}}\nolimits}
\def\NL{\mathop{N\!L}\nolimits}
\def\CH{\mathop{\mathrm{CH}}\nolimits}
\def\proetale{{pro\text{-}\acute{e}tale}}
\def\etale{{\acute{e}tale}}
\def\QCoh{\mathit{QCoh}}
\def\Ker{\mathop{\mathrm{Ker}}}
\def\Im{\mathop{\mathrm{Im}}}
\def\Coker{\mathop{\mathrm{Coker}}}
\def\Coim{\mathop{\mathrm{Coim}}}

% Boxtimes
%
\DeclareMathSymbol{\boxtimes}{\mathbin}{AMSa}{"02}

%
% Macros for moduli stacks/spaces
%
\def\QCohstack{\mathcal{QC}\!\mathit{oh}}
\def\Cohstack{\mathcal{C}\!\mathit{oh}}
\def\Spacesstack{\mathcal{S}\!\mathit{paces}}
\def\Quotfunctor{\mathrm{Quot}}
\def\Hilbfunctor{\mathrm{Hilb}}
\def\Curvesstack{\mathcal{C}\!\mathit{urves}}
\def\Polarizedstack{\mathcal{P}\!\mathit{olarized}}
\def\Complexesstack{\mathcal{C}\!\mathit{omplexes}}
% \Pic is the operator that assigns to X its picard group, usage \Pic(X)
% \Picardstack_{X/B} denotes the Picard stack of X over B
% \Picardfunctor_{X/B} denotes the Picard functor of X over B
\def\Pic{\mathop{\mathrm{Pic}}\nolimits}
\def\Picardstack{\mathcal{P}\!\mathit{ic}}
\def\Picardfunctor{\mathrm{Pic}}
\def\Deformationcategory{\mathcal{D}\!\mathit{ef}}
\usepackage{cleveref}
\crefname{originaldefinition}{original definition}{original definitions}
\Crefname{originaldefinition}{Original definition}{Original definitions}
\makeatletter
\newlabel{algebra-definition-etale}{{00U1}{source}{}{sourcepointer-1}{}}
\newlabel{algebra-definition-etale@cref}{{[originaldefinition][0][]00U1}{[1][0][]source}{}{}{}}
\makeatother
\crefname{originallemma}{original lemma}{original lemmas}
\Crefname{originallemma}{Original lemma}{Original lemmas}
\makeatletter
\newlabel{algebra-lemma-NAK}{{00DV}{source}{}{sourcepointer-2}{}}
\newlabel{algebra-lemma-NAK@cref}{{[originallemma][0][]00DV}{[1][0][]source}{}{}{}}
\makeatother
\crefname{originallemma}{original lemma}{original lemmas}
\Crefname{originallemma}{Original lemma}{Original lemmas}
\makeatletter
\newlabel{algebra-lemma-completion-flat}{{00MB}{source}{}{sourcepointer-3}{}}
\newlabel{algebra-lemma-completion-flat@cref}{{[originallemma][0][]00MB}{[1][0][]source}{}{}{}}
\makeatother
\crefname{originalreference}{original reference}{original references}
\Crefname{originalreference}{Original reference}{Original references}
\makeatletter
\newlabel{lemma-NL-is-limit}{{0AJS}{source}{}{sourcepointer-4}{}}
\newlabel{lemma-NL-is-limit@cref}{{[originalreference][0][]0AJS}{[1][0][]source}{}{}{}}
\makeatother
\crefname{originalreference}{original reference}{original references}
\Crefname{originalreference}{Original reference}{Original references}
\makeatletter
\newlabel{lemma-equivalent-with-artin}{{0AJU}{source}{}{sourcepointer-5}{}}
\newlabel{lemma-equivalent-with-artin@cref}{{[originalreference][0][]0AJU}{[1][0][]source}{}{}{}}
\makeatother
\title{Etale off V(I) is equivalent to the rig-etale condition on the completion}
\author{The Stacks Project Authors}
\date{stacks-project:restricted.tex:local-snapshot-20261012; GNU FDL 1.2 or later}
\begin{document}
\maketitle
\noindent\textit{Editorial note.} All delivered mathematical text is the author's own, copied from the pinned Stacks Project chapter file restricted.tex; no statement or proof step was rewritten, reordered or removed.\par\medskip
\noindent\textit{Editorial note.} The counted proof is the author's complete argument: the proof proves direction (2) by the annihilator-plus-Nakayama argument producing g in 1 + IB and direction (1) by derived completion of the naive cotangent complex quoted from lemma-NL-is-limit, with the completion flatness fact quoted from Algebra\par\medskip
\noindent\textit{Editorial note.} The label lemma-rig-etale names the counted statement, and the proof environment immediately following it in the source is delivered verbatim as its proof.\par\medskip
\noindent\textit{Editorial note.} The counted statement is carried by the author's own lemma environment in the chapter source.\par\medskip
\noindent\textit{Allowed standard prerequisites (editorial).}\begin{itemize}
\item the equivalent conditions of lemma-equivalent-with-artin (same chapter) for the completed algebra B\textasciicircum{}wedge
\item the naive cotangent complex NL\_\{B/A\} of a finite type algebra has finite cohomology modules, and I-adic completion is exact on finite modules (Algebra, Lemma algebra-lemma-completion-flat)
\item NL\textasciicircum{}wedge\_\{B\textasciicircum{}wedge/A\} is the completion (derived completion) of NL\_\{B/A\} (lemma-NL-is-limit of the same chapter)
\item Nakayama's lemma producing an annihilator of the form g = 1 + x with x in IB (Algebra, Lemma algebra-lemma-NAK)
\item the etale criterion for ring maps via the naive cotangent complex (Algebra, Definition algebra-definition-etale)
\end{itemize}\medskip
\ifdefined\hypertarget\hypertarget{sourcepointer-1}{}\fi
\noindent\textit{Source reference.} Original definition 00U1: sources/stacks-chow/tags.txt, source lines 1095--1095. Prerequisite pointer invoked at proof line(s) 2013; it names a standard fact of a sibling chapter that the counted argument uses, and every step of the counted proof built on it is retained verbatim in the delivered body.\par\medskip
\ifdefined\hypertarget\hypertarget{sourcepointer-2}{}\fi
\noindent\textit{Source reference.} Original lemma 00DV: sources/stacks-chow/tags.txt, source lines 512--512. Prerequisite pointer invoked at proof line(s) 2006; it names a standard fact of a sibling chapter that the counted argument uses, and every step of the counted proof built on it is retained verbatim in the delivered body.\par\medskip
\ifdefined\hypertarget\hypertarget{sourcepointer-3}{}\fi
\noindent\textit{Source reference.} Original lemma 00MB: sources/stacks-chow/tags.txt, source lines 817--817. Prerequisite pointer invoked at proof line(s) 2001; it names a standard fact of a sibling chapter that the counted argument uses, and every step of the counted proof built on it is retained verbatim in the delivered body.\par\medskip
\ifdefined\hypertarget\hypertarget{sourcepointer-4}{}\fi
\noindent\textit{Source reference.} Original reference 0AJS: sources/stacks-restricted/pinned-original-source.tex, source lines 596--596. Prerequisite pointer invoked at proof line(s) 2020; it names a standard fact of the same chapter that the counted argument uses, and every step of the counted proof built on it is retained verbatim in the delivered body.\par\medskip
\ifdefined\hypertarget\hypertarget{sourcepointer-5}{}\fi
\noindent\textit{Source reference.} Original reference 0AJU: sources/stacks-restricted/pinned-original-source.tex, source lines 1810--1810. Prerequisite pointer invoked at proof line(s) 1987; it names a standard fact of the same chapter that the counted argument uses, and every step of the counted proof built on it is retained verbatim in the delivered body.\par\medskip
\begin{lemma}
\label{lemma-rig-etale}
Let $A$ be a Noetherian ring and let $I$ be an ideal.
Let $B$ be a finite type $A$-algebra.
\begin{enumerate}
\item If $\Spec(B) \to \Spec(A)$ is \'etale over $\Spec(A) \setminus V(I)$,
then $B^\wedge$ satisfies the equivalent conditions of
Lemma \ref{lemma-equivalent-with-artin}.
\item If $B^\wedge$ satisfies the equivalent conditions of
Lemma \ref{lemma-equivalent-with-artin},
then there exists $g \in 1 + IB$ such that $\Spec(B_g)$ is \'etale
over $\Spec(A) \setminus V(I)$.
\end{enumerate}
\end{lemma}

\begin{proof}
Assume $B^\wedge$ satisfies the equivalent conditions of
Lemma \ref{lemma-equivalent-with-artin}.
The naive cotangent complex $\NL_{B/A}$ is a complex of finite type
$B$-modules and hence $H^{-1}$ and $H^0$ are finite $B$-modules.
Completion is an exact functor on finite $B$-modules (Algebra,
Lemma \ref{algebra-lemma-completion-flat}) and $\NL_{B^\wedge/A}^\wedge$
is the completion of the complex $\NL_{B/A}$ (this is easy to see
by choosing presentations).
Hence the assumption implies there exists a $c \geq 0$ such that
$H^{-1}/I^nH^{-1}$ and $H^0/I^nH^0$ are annihilated by $I^c$
for all $n$. By Nakayama's lemma (Algebra, Lemma \ref{algebra-lemma-NAK})
this means that $I^cH^{-1}$ and $I^cH^0$ are annihilated by an element
of the form $g = 1 + x$ with $x \in IB$. After inverting $g$
(which does not change the quotients $B/I^nB$)
we see that $\NL_{B/A}$ has cohomology annihilated by $I^c$. Thus
$A \to B$ is \'etale at any prime of $B$ not lying over $V(I)$
by the definition of \'etale ring maps, see
Algebra, Definition \ref{algebra-definition-etale}.

\medskip\noindent
Conversely, assume that $\Spec(B) \to \Spec(A)$ is \'etale over
$\Spec(A) \setminus V(I)$. Then for every $a \in I$ there exists
a $c \geq 0$ such that multiplication by $a^c$ is zero $\NL_{B/A}$.
Since $\NL_{B^\wedge/A}^\wedge$ is the derived completion of
$\NL_{B/A}$ (see Lemma \ref{lemma-NL-is-limit}) it follows that
$B^\wedge$ satisfies the equivalent conditions of
Lemma \ref{lemma-equivalent-with-artin}.
\end{proof}

\end{document}
